\chapter{Эмбриональный граф: с чего начинается память}
\label{ch:embryonic}

Пустой граф памяти бесполезен: у существа не бывает мышления «с нуля» — есть
врождённый здравый смысл, на который ложится личный опыт. Поэтому память Рин
начинается с \emph{эмбрионального графа} — нейтрального по характеру каркаса из
готовых баз здравого смысла, который затем обрастает личным опытом через fusion и
консолидацию. Эта глава отвечает предметно: из чего он собран, как выглядят его
узлы, как их сшить и как подключить эмбеддинги.

\section{Две базы, две разные роли}

Нужны две дополняющие базы (Wikidata как содержание не берём — человек не помнит
все столицы; глава~\ref{ch:substrate}):

\begin{itemize}
\item \textbf{ConceptNet} — \emph{предметный и таксономический} здравый смысл: что
чем является, из чего состоит, для чего используется. Даёт Рин «физику мира».
\item \textbf{ATOMIC} — \emph{причинно-социальный} здравый смысл о событиях и
ментальных состояниях: чего человек хотел, что почувствовал, что за этим следует.
Даёт Рин «психологию».
\end{itemize}
Порознь одного мало: ConceptNet без ATOMIC — мир без психологии; ATOMIC без
ConceptNet — реакции без предметного мира.

\section{Как выглядят узлы и рёбра ConceptNet}

ConceptNet — это уже граф. Узел — концепт (слово или короткая фраза), заданный
URI; ребро — типизированная взвешенная ассоциация. Схематически запись:

\begin{lstlisting}[style=pseudo]
узлы:   /c/ru/кружка        /c/ru/посуда      /c/ru/пить      /c/ru/горячий
рёбра (concept, relation, concept, weight):
   (/c/ru/кружка, IsA,       /c/ru/посуда,  2.5)
   (/c/ru/кружка, UsedFor,   /c/ru/пить,    1.8)
   (/c/ru/кружка, HasProperty,/c/ru/горячий,1.2)
\end{lstlisting}

Типов отношений около трёх-четырёх десятков (\textsc{IsA}, \textsc{PartOf},
\textsc{UsedFor}, \textsc{HasProperty}, \textsc{RelatedTo}, \textsc{Causes},
\textsc{MotivatedByGoal}, \dots), и рёбра многоязычны (связывают понятия между
языками). \textbf{Важно для нас}: вес ребра $\approx$ мягкая уверенность из коробки,
а типы отношений прямо ложатся на нашу типизацию рёбер (глава~\ref{ch:substrate}).

\section{Как выглядят узлы ATOMIC}

ATOMIC устроен иначе, и это источник главной сложности сшивки. Его узлы — не
концепты, а \emph{события-фразы} с плейсхолдерами действующих лиц
(\textsc{PersonX}, \textsc{PersonY}). От события идут рёбра фиксированных
\emph{социально-причинных} типов. Схематически:

\begin{lstlisting}[style=pseudo]
событие: "PersonX breaks a cup"
   xIntent  -> (none / accident)     # чего PersonX хотел
   xReact   -> (sad, upset)          # что PersonX почувствовал
   xEffect  -> (cleans it up, buys a new one)   # что произошло с PersonX
   xNeed    -> (to be holding a cup) # что требовалось до
   oReact   -> (concerned)           # что почувствовали окружающие
\end{lstlisting}

То есть «кружки» как отдельного узла в ATOMIC \textbf{нет} — есть событие «разбить
кружку» и его социально-причинные следствия. ATOMIC10x — та же структура, но в
несколько миллионов примеров, порождённых языковой моделью и отфильтрованных
критиком (больше покрытия, больше шума).

\section{Сшивка: главная инженерная работа}

Соблазн сказать «форматы отличаются только на диске» — недооценка. Внутри графовой
базы это одна структура (узлы-токены, направленные типизированные рёбра), и обход от
типа ребра не меняется. Но чтобы Рин, услышав «разбил кружку», прошла от
\emph{события} ATOMIC к \emph{предметным} свойствам кружки из ConceptNet, события и
концепты надо \emph{связать}. Это entity/event linking — нетривиальный шаг.

\subsection{Мост через контент-слова и Numberbatch}

Алгоритм сшивки: из события-фразы ATOMIC извлекаем контентные слова (существительные,
глаголы) и линкуем их к концепт-узлам ConceptNet — сначала точным лемма-матчем,
затем эмбеддинг-фолбэком.

\begin{pitfall}[Не bag-of-words и не one-hot]
Грубые представления (мешок слов, one-hot) не знают синонимии: «mug» и «cup»,
«разбить» и «сломать» окажутся разными токенами, и мост выйдет дырявым. Нужен
эмбеддинг, знающий смысловую близость.
\end{pitfall}

\emph{ConceptNet Numberbatch} — родной эмбеддинг самого ConceptNet: каждому
концепт-узлу сопоставлен вектор ($\sim$300 измерений), причём смыслово близкие
концепты близки в пространстве, а языки согласованы. Это идеальный мост:

\begin{lstlisting}[style=pseudo]
def link_atomic_event_to_concepts(event_phrase, conceptnet, numberbatch):
    concepts = []
    for token in content_words(event_phrase):     # "break", "cup"
        node = conceptnet.lemma_match(token)       # 1) точный матч
        if node is None:                           # 2) эмбеддинг-фолбэк
            node = numberbatch.nearest(embed(token))
        concepts.append(node)
    return concepts     # событие ATOMIC <-> концепты ConceptNet
\end{lstlisting}

Тот же Numberbatch решает ещё две задачи (одно пространство — три применения):
член релевантности $\text{rel}(B\mid q)$ в правиле обхода (глава~\ref{ch:walker}) и
сидирование валентности узлов (регрессия аффект-лексикона в это же пространство).

\subsection{Заземление действующих лиц}

Плейсхолдеры \textsc{PersonX}/\textsc{PersonY} — не абстракции, а слоты. В рантайме
они \emph{заземляются} на конкретных людей: \textsc{PersonX}~$\mapsto$ зритель $Z$,
\textsc{PersonY}~$\mapsto$ второй пользователь проекта. Так общий социальный шаблон
превращается в запись о конкретном отношении.

\begin{workedbox}[«Разбил кружку» на эмбриональном графе Рин]
Зритель пишет «разбил любимую кружку». (1)~Событие матчится к ATOMIC-шаблону
«PersonX breaks a cup», \textsc{PersonX}~$\mapsto Z$. (2)~Через ATOMIC всплывают
социально-причинные узлы: досада ($Z$), сочувствие (окружающие), «нужна новая».
(3)~Через сшивку с ConceptNet всплывают предметные свойства: кружка — посуда, из
неё пьют, бывает любимой. (4)~Если в личной памяти Рин есть узел «$Z$ давно хотел
новую кружку» ($-$ contradicts «досада»), появляется основание для уместного
«поздравляю» вместо сочувствия. Ни одна база в одиночку этого не даёт: ATOMIC —
эмоцию, ConceptNet — предмет, личная память — контекст.
\end{workedbox}

\section{Как строить эмбрион: засеянный подграф, а не слепой импорт}
\label{sec:embryo-build}

Здесь важная инженерная развилка. \emph{Слепо залить} весь ConceptNet и весь ATOMIC
в память — плохо: получишь много шума (банальные и устаревшие связи, англоцентризм,
странности social commonsense, слабый перенос на русский, низкую релевантность
именно домашней Рин). Тогда ты будешь отлаживать не память Рин, а шумы чужой базы
здравого смысла. Но и другая крайность — \emph{писать граф руками} на 500 узлов —
тоже плоха: медленно, и такой граф слишком разрежён, чтобы на нём проявились
топологические эффекты (хабы, боттлнеки, over-squashing), которые как раз надо
тестировать (глава~\ref{ch:topology}).

Правильный — \textbf{серединный путь}: извлечь \emph{небольшой засеянный проверяемый
подграф}.
\begin{enumerate}\itemsep1pt
\item \textbf{Сид} — концепты реального мира Рин: дом, код, усталость, кошка, универ,
дота, DTI, RIN, бытовые действия, режимы общения.
\item \textbf{Расширение} — ego-граф на $k$ хопов от сида по ConceptNet/ATOMIC.
\item \textbf{Фильтр по весу ребра} — отсекает банальщину и слабые ассоциации
(вес ConceptNet $\approx$ мягкая уверенность из коробки).
\item \textbf{Ручная чистка глазами} — один раз, для пары тысяч узлов дёшево.
\item \textbf{Статус} — это \emph{источник кандидатов}, а не замороженная истина:
дальше Рин его сама растит и прунит через использование.
\end{enumerate}
Так получаем реальную структуру (для топологических тестов) $+$ реальные ассоциации
(не рукописные) $+$ проверяемость глазами $+$ релевантность (засеяно на её домены), и
обходим шум слепого импорта.

\begin{openq}[Русский ConceptNet тоньше английского]
Практический кавеат: русскоязычная часть ConceptNet заметно беднее английской.
Варианты: держать \emph{концептный backbone} в англо-пространстве (там связей
больше), а русскую поверхность накручивать сверху; либо принять более разрежённый
русский и дополнять его вручную/из других источников. Что даст лучшее качество на
домашних сценариях Рин — открыто, проверяется на стенде (глава~\ref{ch:experiment}).
\end{openq}

\section{Как эмбриональный граф превращается в память Рин}

Эмбриональный граф \emph{нейтрален по характеру}: это «то, что досталось от
предков-разработчиков», усреднённый здравый смысл. Личность и личный опыт входят не
переписыванием этих узлов, а \emph{надстройкой}: каждое общение порождает через
fusion (глава~\ref{ch:fusion}) новые узлы и рёбра с временными метками, а
консолидация закрепляет устойчивое. Со временем поверх нейтрального каркаса
вырастает индивидуальная, окрашенная память — но фундамент здравого смысла остаётся
как \emph{система координат}, относительно которой отклонения читаются как характер
(глава~\ref{ch:personality}).

\begin{intuition}[Эмбрион, а не энциклопедия]
Название «эмбриональный» выбрано не случайно. Это не готовое знание о мире, которое
Рин «зачитывает», а зачаток, минимальный каркас связей, на котором \emph{развивается}
собственная память. Врождённые узлы со временем могут и потускнеть, и быть
пересмотрены (континуум глубины консолидации, глава~\ref{ch:substrate}) — просто
медленнее приобретённых.
\end{intuition}
